% -------------------------------------------------------------------------

    \subsection{Теоретический вопрос III.1: Выведите characteristic equation for eigenvalues}
\textbf{Формулировка вопроса:} Выведите характеристическое уравнение для собственных значений матрицы из определения
собственного вектора.

\textbf{Ответ и строгий вывод:}

По определению ненулевой вектор $v \in \mathbb{C}^n$ ($v \ne 0$) называется собственным вектором матрицы $A \in
\mathbb{R}^{n \times n}$, отвечающим собственному значению $\lambda \in \mathbb{C}$, если: \[ A v = \lambda v.
\] Перенесем все слагаемые в левую часть, используя единичную матрицу $I_n$: \[ A v - \lambda I_n v = 0 \iff (A -
\lambda I_n) v = 0.
\] Мы получили однородную систему линейных уравнений относительно вектора $v$ с квадратной матрицей коэффициентов
$M(\lambda) \coloneqq A - \lambda I_n \in \mathbb{C}^{n \times n}$.

Условие существования \textbf{ненулевого} собственного вектора $v \ne 0$: Однородная система $M(\lambda) v = 0$ имеет
нетривиальное решение тогда и только тогда, когда ядро матрицы $M(\lambda)$ содержит ненулевые векторы: \[ \ker(A -
\lambda I_n) \ne \{0\} \iff \nullity(A - \lambda I_n) > 0.
\] По теореме о ранге и дефекте это равносильно тому, что ранг матрицы строго меньше ее размера: \[ \rank(A - \lambda
I_n) < n.
\] А квадратная матрица неполного ранга является \textbf{вырожденной}, что по критерию обратимости эквивалентно
равенству нулю ее определителя: \[ \det(A - \lambda I_n) = 0.
\] Многочлен $p_A(\lambda) \coloneqq \det(A - \lambda I_n)$ является многочленом степени $n$ относительно переменной
$\lambda$ и называется \textbf{характеристическим многочленом}, а уравнение $\det(A - \lambda I_n) = 0$ —
\textbf{характеристическим уравнением}.

% -------------------------------------------------------------------------

    \subsection{Теоретический вопрос III.2: Почему distinct eigenvalues дают independent eigenvectors?}
\textbf{Формулировка вопроса:} Докажите теорему: собственные векторы, отвечающие попарно различным собственным
значениям, линейно независимы.

\textbf{Ответ и строгое доказательство:}

    \begin{theorem}
Пусть $v_1, v_2, \dots, v_k$ — собственные векторы матрицы $A$, отвечающие попарно различным собственным значениям
$\lambda_1, \lambda_2, \dots, \lambda_k$ ($\lambda_i \ne \lambda_j$ при $i \ne j$).
Тогда система векторов $\{v_1, \dots, v_k\}$ линейно независима.
    \end{theorem}
    \begin{proof}
Доказательство проведем методом математической индукции по числу векторов $k$.

\textbf{База индукции ($k = 1$):} По определению собственный вектор ненулевой: $v_1 \ne 0$.
Один ненулевой вектор всегда образует линейно независимую систему ($c_1 v_1 = 0 \implies c_1 = 0$).
База доказана.

\textbf{Индукционный переход:} Пусть утверждение справедливо для $k - 1$ векторов: любые $k - 1$ собственных векторов с
различными собственными значениями линейно независимы.
Рассмотрим систему из $k$ векторов и составим их линейную комбинацию, равную нулю:
        \begin{equation} \label{eq:t32_1} c_1 v_1 + c_2 v_2 + \dots + c_{k-1} v_{k-1} + c_k v_k = 0. \end{equation}
Применим оператор $A$ к обеим частям равенства \eqref{eq:t32_1}.
По линейности и свойству $A v_i=\lambda_i v_i$:
        \begin{equation} \label{eq:t32_2}
            c_1 \lambda_1 v_1 + c_2 \lambda_2 v_2 + \dots + c_{k-1} \lambda_{k-1} v_{k-1} + c_k \lambda_k v_k = 0.
        \end{equation}
Теперь умножим исходное уравнение \eqref{eq:t32_1} на скаляр $\lambda_k$:
        \begin{equation} \label{eq:t32_3}
            c_1 \lambda_k v_1 + c_2 \lambda_k v_2 + \dots + c_{k-1} \lambda_k v_{k-1} + c_k \lambda_k v_k = 0.
        \end{equation}
Вычтем почленно равенство \eqref{eq:t32_3} из равенства \eqref{eq:t32_2}: \[ c_1 (\lambda_1 - \lambda_k) v_1 + c_2
(\lambda_2 - \lambda_k) v_2 + \dots + c_{k-1} (\lambda_{k-1} - \lambda_k) v_{k-1} + 0 = 0.
\] Мы получили равенство нулю линейной комбинации первых $k - 1$ векторов $\{v_1, \dots, v_{k-1}\}$.
По индукционному предположению эта система линейно независима.
Следовательно, все ее коэффициенты обязаны быть равны нулю: \[ c_i (\lambda_i - \lambda_k) = 0 \quad \text{для всех } i
\in \{1, 2, \dots, k-1\}.
\] Так как по условию все собственные значения попарно различны, разность $\lambda_i - \lambda_k \ne 0$ при $i< k$.
Отсюда получаем: \[ c_1 = c_2 = \dots = c_{k-1} = 0.
\] Подставим найденные нулевые коэффициенты обратно в исходное равенство \eqref{eq:t32_1}: \[ 0 + 0 + \dots + 0 + c_k
v_k = 0 \iff c_k v_k = 0.
\] Поскольку собственный вектор $v_k \ne 0$, получаем $c_k = 0$.
Итак, все коэффициенты $c_1 = c_2 = \dots = c_k = 0$.
Система векторов $\{v_1, \dots, v_k\}$ линейно независима, что завершает шаг индукции.
    \end{proof}

% -------------------------------------------------------------------------

    \subsection{Теоретический вопрос III.3: Объясните diagonalization through eigenbasis}
\textbf{Формулировка вопроса:} Объясните процедуру и геометрический смысл диагонализации матрицы через базис из
собственных векторов.

\textbf{Ответ:}

Пусть квадратная матрица $A \in \mathbb{R}^{n \times n}$ обладает $n$ линейно независимыми собственными векторами $v_1,
v_2, \dots, v_n$, отвечающими собственным значениям $\lambda_1, \dots, \lambda_n$: \[ A v_j = \lambda_j v_j, \quad j =
1, \dots, n.
\] Составим матрицу $P \in \mathbb{R}^{n \times n}$, столбцами которой являются эти собственные векторы: \[ P =
\begin{pmatrix} | & | & & | \\ v_1 & v_2 & \dots & v_n \\ | & | & & |
    \end{pmatrix}. \]
Поскольку столбцы матрицы $P$ линейно независимы, матрица $P$ невырождена ($\det P \ne 0$) и существует обратная матрица
$P^{-1}$.

Умножим матрицу $A$ на матрицу $P$ по правилу блочного умножения: \[ A P = A \begin{pmatrix} | & & | \\ v_1 & \dots &
v_n \\ | & & |
    \end{pmatrix} = \begin{pmatrix}
| & & | \\ A v_1 & \dots & A v_n \\ | & & |
    \end{pmatrix} = \begin{pmatrix}
| & & | \\ \lambda_1 v_1 & \dots & \lambda_n v_n \\ | & & |
    \end{pmatrix}. \]
С другой стороны, умножим матрицу $P$ справа на диагональную матрицу $\Lambda = \diag(\lambda_1, \dots, \lambda_n)$: \[
P \Lambda = \begin{pmatrix} | & & | \\ v_1 & \dots & v_n \\ | & & |
    \end{pmatrix} \begin{pmatrix}
\lambda_1 & & 0 \\ & \ddots & \\ 0 & & \lambda_n
    \end{pmatrix} = \begin{pmatrix}
| & & | \\ \lambda_1 v_1 & \dots & \lambda_n v_n \\ | & & |
    \end{pmatrix}. \]
Столбцы обеих матриц тождественно равны, следовательно: \[ A P = P \Lambda.
\] Умножив равенство справа на $P^{-1}$, получаем формулу \textbf{диагонализации}: \[ A = P \Lambda P^{-1}.
\] А умножив слева на $P^{-1}$, получаем формулу перехода к диагональному виду: \[ \Lambda = P^{-1} A P. \]

\textbf{Геометрический смысл диагонализации:} Матрица $P$ является матрицей перехода от стандартного базиса к
собственному базису $\mathcal{B}=(v_1, \dots, v_n)$.
Преобразование $y = Ax$ в собственном базисе распадается на $n$ независимых одномерных масштабирований: \[
[y]_\mathcal{B} = \Lambda [x]_\mathcal{B} \iff \begin{pmatrix} y'_1 \\ \vdots \\ y'_n
    \end{pmatrix} = \begin{pmatrix}
\lambda_1 x'_1 \\ \vdots \\ \lambda_n x'_n
    \end{pmatrix}. \]
Собственные направления — это оси естественной системы координат оператора, вдоль которых отсутствует перемешивание
координат.

    \newpage